\chapter{Détails de validation et reproductibilité}
\label{annexe:validation}

\newcommand{\annexeheading}[1]{%
  \par\medskip
  \noindent\textbf{#1}\par\smallskip
}

% =====================================================================
\annexeheading{Paramètres figés}
% =====================================================================

La configuration finale retenue est gelée après les campagnes de validation. Elle fixe
notamment un intervalle de 15~minutes, une contamination de 0.06, une fenêtre de persistance
de 7~heures, un minimum de 2~bins anomaux par fenêtre, un mode de cohérence multi-signaux
(MIX) avec un seuil de taux d'anomalies de 0.24, un cooldown de 12~heures et un seuil
z-score sur le \textit{Motion Index} de 2.2. Le Tableau~\ref{tab:config_finale}
(Chapitre~\ref{chapitre:methodologie}) présente ces paramètres en détail. Un fichier de
référence documente la provenance expérimentale de ces choix et les campagnes de validation
associées.

% =====================================================================
\annexeheading{Grille de paramètres explorée}
% =====================================================================

Le Tableau~\ref{tab:grille_complete} présente la grille complète de paramètres explorée au
cours des campagnes de validation robuste. Cette grille a été élargie progressivement de
v4 (12~configurations) à v5.2 (336~configurations).

\begin{table}[H]
    \centering
    \small
    \caption{Grille de paramètres explorée dans les validations robustes. Les colonnes
    \texttt{v4}/\texttt{v5.1}/\texttt{v5.2} donnent le nombre de valeurs testées pour
    chaque paramètre~; le nombre de configurations est le \emph{produit} de ces comptes.}
    \label{tab:grille_complete}
    \setlength{\tabcolsep}{4pt}
    \begin{tabular}{llccc}
        \toprule
        \textbf{Paramètre} & \textbf{Valeurs explorées (union des campagnes)}
            & \textbf{v4} & \textbf{v5.1} & \textbf{v5.2} \\
        \midrule
        \texttt{contamination}    & 0.05, 0.06, 0.07, 0.08
            & 3 & 3 & 4 \\
        \texttt{persist\_hours}   & 6, 7, 8
            & 2 & 2 & 3 \\
        \texttt{alert\_min}       & 2, 3
            & 1 & 1 & 2 \\
        \texttt{mix\_mode}        & MIX
            & 1 & 1 & 1 \\
        \texttt{mix\_rate\_thr}   & 0.16, 0.18, 0.20, 0.22, 0.24, 0.28, 0.34
            & 2 & 4 & 7 \\
        \texttt{mi\_z\_high\_thr} & 2.2, 2.5
            & 1 & 2 & 2 \\
        \midrule
        \textbf{Total configs} ($=$ produit) &
            & \textbf{12} & \textbf{48} & \textbf{336} \\
        \textbf{Seeds}            &
            & 8 & 8 & 8 \\
        \textbf{Scénarios}        &
            & 3 & 3 & 3 \\
        \textbf{Total exécutions}       &
            & \textbf{288} & \textbf{1\,152} & \textbf{8\,064} \\
        \bottomrule
    \end{tabular}
\end{table}

L'élargissement progressif de la grille permet de vérifier que les seuils retenus ne sont
pas un optimum local fragile, mais une zone de stabilité dans l'espace des paramètres. La
colonne ``v5.2'' correspond à la grille la plus large, utilisée pour le verdict final de
robustesse.

% =====================================================================
\annexeheading{Campagnes de référence}
% =====================================================================

Les principales campagnes de validation utilisées dans ce mémoire sont les suivantes:
\begin{itemize}
  \item \textit{Validation manuelle v2}: première campagne d'injections détectables et
  non-détectables, servant de référence historique;
  \item \textit{Validation manuelle v2.1}: campagne de consolidation alignée sur les seuils
  finaux (\texttt{persist\_hours} = 7);
  \item \textit{Validation robuste v4}: 288~exécutions, grille de 12~configurations,
  3~scénarios, 8~seeds;
  \item \textit{Validation robuste v5.1}: 1\,152~exécutions, grille élargie à 48~configurations,
  ajout des paramètres \texttt{mi\_z\_high\_thr} et de seuils supplémentaires;
  \item \textit{Validation robuste v5.2}: 8\,064~exécutions, grille complète de 336~configurations,
  vérification de stabilité à grande échelle;
  \item \textit{Validation exploratoire LOF}: 288~exécutions selon le protocole v4, avec LOF
  comme détecteur en remplacement d'IF (comparaison au niveau événement);
  \item \textit{Benchmark de performance}: mesure de scalabilité sur le corpus complet,
  avec décomposition du temps par composante;
  \item \textit{Ablation étendue}: 90~cas par variante (3~vaches $\times$ 3~profils
  $\times$ 10~seeds), comparaison IF/LOF/règles seules.
\end{itemize}

Au total, les campagnes de validation représentent 9\,504~exécutions expérimentales
indépendantes, générant plus de 28\,000~évaluations événementielles individuelles.

% =====================================================================
\annexeheading{Détails des profils d'injection}
% =====================================================================

Les injections synthétiques modifient les colonnes brutes du capteur (\textit{Steps}, \textit{Motion Index},
\textit{Lying Time}, \textit{Standing Time}, \textit{Transitions}) en appliquant des multiplicateurs par phases.
Le Tableau~\ref{tab:injection_multiplicateurs} détaille les amplitudes utilisées pour le
profil \texttt{detectable\_strong}~; les phases y alternent entre repos (proche de 0)
et épisode (valeurs élevées) toutes les 6~bins ($\approx$1\,h\,30).

\begin{table}[H]
    \centering
    \caption{Multiplicateurs d'injection pour le profil \texttt{detectable\_strong}.}
    \label{tab:injection_multiplicateurs}
    \begin{tabular}{lcc}
        \toprule
        \textbf{Signal}       & \textbf{Phase repos}   & \textbf{Phase épisode} \\
        \midrule
        \textit{Steps}                 & 0.02 -- 0.10           & 1.80 -- 2.80 \\
        \textit{Motion Index}          & 0.01 -- 0.08           & 2.00 -- 3.40 \\
        \textit{Lying Time}            & 2.60 -- 4.20           & 0.65 -- 1.10 \\
        \textit{Standing Time}         & 0.04 -- 0.20           & 1.25 -- 2.10 \\
        \textit{Transitions}           & 0.03 -- 0.15           & 3.00 -- 5.20 \\
        \bottomrule
    \end{tabular}
\end{table}

Les multiplicateurs sont tirés uniformément dans les intervalles indiqués, ce qui introduit
une variabilité réaliste entre injections. Le Tableau~\ref{tab:injection_tous_profils}
compare les quatre profils d'injection utilisés dans les campagnes de validation~; les
valeurs indiquent l'intervalle [min, max] des multiplicateurs appliqués aux signaux
bruts.

\begin{table}[H]
    \centering
    \caption{Multiplicateurs d'injection pour les quatre profils.}
    \label{tab:injection_tous_profils}
    \small
    \begin{tabular}{lcccc}
        \toprule
        \textbf{Signal} & \textbf{Strong} & \textbf{Multisignal} & \textbf{Borderline}
            & \textbf{Hidden} \\
        \midrule
        \textit{Steps} (épisode)     & 1.80--2.80 & 1.40--2.20 & 1.35--1.95 & $\sim$0.86 \\
        \textit{Motion Index} (ép.)  & 2.00--3.40 & 1.60--2.80 & 1.70--2.45 & $\sim$0.86 \\
        \textit{Lying Time} (ép.)    & 0.65--1.10 & 0.70--1.30 & 0.62--0.98 & 1.00 \\
        \textit{Standing Time} (ép.) & 1.25--2.10 & 1.10--1.90 & 1.15--1.70 & 1.00 \\
        \textit{Transitions} (ép.)   & 3.00--5.20 & 2.50--4.80 & 1.85--3.05 & $\sim$0.86 \\
        \midrule
        Familles touchées   & 5/5        & 5/5        & 5/5        & 1/5 \\
        Amplitude relative  & forte      & modérée    & faible     & quasi-nulle \\
        \bottomrule
    \end{tabular}
\end{table}

Le profil \textit{hidden} (non détectable) constitue le contrôle négatif: il ne modifie
qu'une seule famille de signaux avec un multiplicateur proche de~1
($\mathcal{N}(0.86, 0.04)$, clippé à $[0.70, 0.98]$), ce qui produit une variation
indistinguable du bruit naturel. La graduation des amplitudes entre \textit{strong},
\textit{multisignal}, borderline et \textit{hidden} permet d'évaluer la
sensibilité du pipeline à différents niveaux de sévérité.

% =====================================================================
\annexeheading{Artefacts produits par campagne}
% =====================================================================

Chaque campagne de validation produit un ensemble standardisé d'artefacts, présentés dans
le Tableau~\ref{tab:artefacts}.

\begin{table}[H]
    \centering
    \caption{Artefacts produits par chaque campagne de validation.}
    \label{tab:artefacts}
    \begin{tabular}{lp{8cm}}
        \toprule
        \textbf{Artefact}                & \textbf{Contenu} \\
        \midrule
        Métadonnées de campagne          & Paramètres, grille, seeds, horodatage,
                                           durée d'exécution \\
        Résumé par configuration         & Métriques agrégées (détection, IoU, faux positifs),
                                           verdict GO/NO-GO, score de sélection \\
        Rapport par scénario             & Métriques détaillées par scénario d'injection
                                           (strong, borderline, hidden) \\
        Évaluation événementielle        & Une ligne par événement injecté, avec chevauchement,
                                           IoU, raison de détection/non-détection \\
        Résultats par exécution                & Paramètres et métriques de chaque exécution individuelle \\
        \bottomrule
    \end{tabular}
\end{table}

% =====================================================================
\annexeheading{Score de sélection v4}
% =====================================================================

Le classement des configurations dans la validation robuste repose sur un score de
sélection pondéré:

\begin{equation}
    S = 0.50 \cdot \overline{\text{IoU20}}_{\text{strong}}
      + 0.20 \cdot \overline{\text{det}}_{\text{strong}}
      + 0.20 \cdot \overline{\text{det}}_{\text{border}}
      - 0.25 \cdot \overline{\text{leak}}_{\text{hidden}}
      - 0.30 \cdot \overline{\text{FN}}
      - 0.10 \cdot \sigma_{\text{FN}}
    \label{eq:selection_score}
\end{equation}

\noindent
où $\overline{\text{IoU20}}_{\text{strong}}$ est le recouvrement strict moyen sur les cas
forts, $\overline{\text{det}}$ les taux de détection, $\overline{\text{leak}}$ le taux de
fuite, $\overline{\text{FN}}$ le taux moyen de fausses notifications et $\sigma_{\text{FN}}$
son écart-type.

La pondération reflète trois priorités opérationnelles hiérarchisées:
\begin{enumerate}
    \item \textit{Précision temporelle} (50\% pour IoU20): pour un éleveur, savoir
    \textit{quand} un épisode a commencé est plus utile qu'une simple alerte binaire.
    Un système qui détecte mais localise mal dans le temps perd sa valeur pratique;
    \item \textit{Maîtrise des fausses alertes} ($-$30\% pour $\overline{\text{FN}}$,
    $-$10\% pour $\sigma_{\text{FN}}$): la littérature montre que les éleveurs cessent
    de consulter les systèmes d'alerte dès que le taux de faux positifs dépasse
    $\sim$0.2/cow-day \citep{rutten2013invited}. La pénalisation de la variance
    ($\sigma$) garantit en outre une fiabilité prédictible;
    \item \textit{Couverture de détection} (20\% chacun pour strong et borderline):
    le poids symétrique entre cas forts et borderline évite de favoriser des configurations
    agressives sur les cas faciles au détriment des cas ambigus.
\end{enumerate}
La pénalité de fuite ($-$25\%) est élevée car une alerte sur un cas volontairement non
détectable signale une défaillance grave de la spécificité. Ces poids ne sont pas issus
d'une optimisation formelle mais d'un arbitrage explicite, documenté et gelé avant
l'exécution des campagnes de validation.

% =====================================================================
\annexeheading{Intégrité des résultats}
% =====================================================================

L'intégrité des artefacts de référence est vérifiée via un manifeste cryptographique
SHA-256. Ce manifeste contient l'empreinte de chaque fichier de résultats produit par les
campagnes de validation. La vérification consiste à recalculer les empreintes et à les
comparer au manifeste de référence.

Le résultat attendu pour la baseline officielle est: zéro fichier manquant, zéro fichier
modifié, zéro fichier inattendu. Cette lecture signifie qu'aucun artefact de référence n'a
disparu, qu'aucun résultat n'a été modifié en place, et qu'aucun bruit non prévu n'a été
injecté dans la baseline vérifiée. Ce mécanisme garantit que les réorganisations de code
n'ont pas altéré les résultats expérimentaux sur lesquels reposent les conclusions du
mémoire.

% =====================================================================
\annexeheading{Archive logicielle du mémoire}
% =====================================================================

L'archive logicielle de référence du mémoire est décrite dans le manifeste
\path{data/memoire_archive_manifest_v2.json}. Ce manifeste liste:
\begin{enumerate}
    \item le dépôt Git privé de référence du projet (accessible sur demande
    auprès de l'auteur);
    \item le tag de référence de la version figée du mémoire (\texttt{memoire-v2});
    \item les sources LaTeX du manuscrit;
    \item les paramètres gelés (\texttt{data/final\_thresholds\_v1.json});
    \item les manifestes d'intégrité SHA-256;
    \item les principaux CSV de résultats cités dans le manuscrit, y compris les statistiques
    appariées IF/LOF exportées dans \path{data/if_lof_stats/}.
\end{enumerate}

Cette archive ne comprend pas les données brutes capteurs, non diffusables pour des raisons
de confidentialité. En revanche, elle fournit tous les éléments nécessaires pour rejouer
la démarche sur un nouveau corpus disposant du même schéma de colonnes.

% =====================================================================
\annexeheading{Reproduction des expériences}
% =====================================================================

Pour reproduire les résultats présentés dans ce mémoire, les étapes suivantes sont
nécessaires:

\begin{enumerate}
    \item \textit{Installation}: installer les dépendances déclarées dans le fichier de
    dépendances du projet (Python~3.10+, pandas, scikit-learn, NumPy);
    \item \textit{Données}: placer le corpus brut (37\,839~lignes, 28~vaches) dans le
    répertoire de données du projet;
    \item \textit{Validation manuelle}: exécuter les notebooks de validation manuelle
    v2 et v2.1;
    \item \textit{Validation robuste}: exécuter les notebooks v4, v5.1 et v5.2 dans
    l'ordre, chacun produisant ses artefacts horodatés;
    \item \textit{Comparaison IF/LOF}: exécuter \texttt{python scripts/compute\_if\_lof\_paired\_stats.py}
    pour régénérer les intervalles de confiance appariés à partir des CSV existants;
    \item \textit{Vérification}: comparer les empreintes SHA-256 des artefacts produits
    avec le manifeste de référence.
\end{enumerate}

Le seed aléatoire (\texttt{random\_state} = 42 pour IF, seeds explicites pour les
validations) garantit la reproductibilité bit-à-bit des résultats sur une même version
de scikit-learn.

% =====================================================================
\annexeheading{Variabilité inter-vache des fausses notifications}
\label{annexe:variabilite_vache}
% =====================================================================

Les métriques agrégées (0.105 fausse notification/vache-jour, 91.7\% d'accord) résument
le comportement moyen du pipeline sur l'ensemble du troupeau, mais masquent la variabilité
naturelle entre individus. Le Tableau~\ref{tab:distribution_par_vache} classe les 28~vaches
en quatre groupes selon leur taux individuel de fausses notifications, mesurés sur
41~jours pour la configuration finale v5.2~; les groupes sont définis par quartile du
taux individuel et les effectifs exacts sont calculés sur la configuration à seed
médian.

\begin{table}[H]
    \centering
    \small
    \caption{Distribution des fausses notifications par vache sur la période d'étude.}
    \label{tab:distribution_par_vache}
    \begin{tabular}{lccccc}
        \toprule
        \textbf{Groupe} & \textbf{$n$ vaches} & \textbf{FN totales} &
        \textbf{FN/vache} & \textbf{FN/vache-jour} & \textbf{Description} \\
        \midrule
        Très stable  & 7 & 0--1   & $\leq$1   & $\leq$0.024 & Signal quasi-nul \\
        Stable       & 7 & 2--3   & 2--3      & 0.05--0.07  & Comportement régulier \\
        Actif        & 7 & 4--7   & 4--7      & 0.10--0.17  & Variabilité modérée \\
        Très actif   & 7 & 8--15  & 8--15     & 0.20--0.37  & Forte variabilité \\
        \midrule
        \textbf{Total} & \textbf{28} & — & $\mathbf{4.6}$ (moy.) & $\mathbf{0.112}$ & \\
        \bottomrule
    \end{tabular}

    \vspace{0.3em}
    \raggedright\footnotesize\textit{Note}: la valeur totale de 0.112~FN/vache-jour est
    la \emph{moyenne des taux individuels} des 28~vaches (configuration finale, seed
    médian). Elle diffère légèrement du taux agrégé de 0.105~FN/vache-jour rapporté au
    Chapitre~\ref{chapitre:validation}, ce dernier étant calculé comme moyenne \emph{par
    exécution} sur la grille robuste v5.1~; les deux mesures sont cohérentes et restent
    sous le seuil critique de la littérature.
\end{table}

\textit{Mécanisme de la variabilité.}
Les vaches du groupe \textit{Très actif} présentent des profils comportementaux naturellement
irréguliers: pics d'activité locomotrice plus fréquents, alternances couché/debout plus
variables, et \textit{Motion Index} plus élevé en dehors des épisodes injectés. Pour ces animaux,
la distribution des scores Isolation~Forest est décalée vers des valeurs plus élevées en
condition normale, ce qui abaisse le seuil effectif de détection. Ce phénomène, bien
documenté dans la littérature sur la détection d'anomalies
\citep{chandola2009anomaly}, rappelle que le score IF est relatif au profil global de la
vache: une anomalie est toujours définie par rapport à la norme individuelle.

\textit{Implication pour le déploiement.}
Deux adaptations opérationnelles découlent de cette analyse:
\begin{enumerate}
    \item \textit{Calibration individuelle}: un ajustement du seuil de contamination
    par vache permettrait de maintenir un taux uniforme de fausses alertes à l'échelle du
    troupeau, plutôt que d'appliquer un seuil global qui surcharge les animaux naturellement
    actifs;
    \item \textit{Pondération des alertes}: le score de confiance pourrait intégrer
    l'historique individuel de l'animal pour différencier une anomalie forte pour une vache
    calme d'une anomalie modérée pour une vache active.
\end{enumerate}

Ces adaptations constituent des axes d'amélioration concrets, validables sur un corpus
élargi. Le taux agrégé (0.105~FN/vache-jour) reste bien en dessous du seuil critique
identifié dans la littérature (${\sim}0.2$~FN/vache-jour, \cite{rutten2013invited}).

% =====================================================================
\annexeheading{Distribution des recouvrements temporels (IoU)}
\label{annexe:distribution_iou}
% =====================================================================

La Figure~\ref{fig:iou_distribution} présente les distributions de best\_iou pour les
trois scénarios d'injection sur l'ensemble des configurations GO de la campagne v5.2.

\begin{figure}[H]
    \centering
    \begin{tikzpicture}
        \begin{axis}[
            ybar,
            bar width=8pt,
            width=12cm, height=7cm,
            ylabel={Fréquence relative (\%)},
            xlabel={Best IoU},
            symbolic x coords={0.00, 0.05, 0.10, 0.15, 0.20, 0.25, 0.30, 0.35, 0.40, 0.45, 0.50+},
            xtick=data,
            x tick label style={font=\scriptsize, rotate=30, anchor=east},
            ymin=0, ymax=60,
            ymajorgrids=true,
            grid style={black!8},
            legend style={at={(0.98,0.96)}, anchor=north east, font=\small,
                          legend columns=1, draw=black!20, fill=white},
            enlarge x limits=0.08,
            axis line style={black!60},
            tick style={black!40},
        ]
            \addplot[fill=rulegreen!40, draw=rulegreen!80] coordinates {
                (0.00,0) (0.05,0) (0.10,2) (0.15,5) (0.20,8)
                (0.25,15) (0.30,22) (0.35,25) (0.40,15) (0.45,5) (0.50+,3)
            };
            \addlegendentry{Strong}

            \addplot[fill=loforange!40, draw=loforange!80] coordinates {
                (0.00,42) (0.05,18) (0.10,15) (0.15,10) (0.20,7)
                (0.25,4) (0.30,2) (0.35,1) (0.40,1) (0.45,0) (0.50+,0)
            };
            \addlegendentry{Borderline}

            \addplot[fill=profgray!30, draw=profgray!60] coordinates {
                (0.00,55) (0.05,0) (0.10,0) (0.15,0) (0.20,0)
                (0.25,0) (0.30,0) (0.35,0) (0.40,0) (0.45,0) (0.50+,0)
            };
            \addlegendentry{Non-détectable}
        \end{axis}
    \end{tikzpicture}
    \caption{Distribution des \texttt{best\_iou} par scénario d'injection (campagne v5.2).}
    \label{fig:iou_distribution}
\end{figure}

La distribution des cas forts est centrée autour de 0.30--0.35, avec une queue s'étendant
jusqu'à 0.50+, ce qui signifie que le pipeline recouvre environ un tiers de la fenêtre
temporelle de l'épisode injecté dans la majorité des configurations stables. La
distribution bimodale des cas borderline, avec un pic à zéro et une queue jusqu'à 0.20,
confirme que ces cas se situent à la frontière de détection du pipeline.
